\documentclass[preview]{standalone}

\usepackage{amsmath}
\usepackage{amssymb}
\usepackage{stellar}
\usepackage{definitions}
\usepackage{bettelini}

\begin{document}

\id{sequences-functions-integration}
\genpage

\section{Interchange with integration}

\begin{snippettheorem}{uniform-convergence-integration-theorem}{Integration of a uniformly convergent sequence}
    Let \(a<b\) and let \(\{f_n\}_{n\in\naturalnumbers}\) be a \sequence of
    Riemann-integrable \function[functions]
    \(f_n\colon[a,b]\fromto\realnumbers\). If \(f_n\) converges uniformly on
    \([a,b]\) to \(f\colon[a,b]\fromto\realnumbers\), then \(f\) is Riemann
    integrable and
    \[
        \integral[a][b][f(x)][x]
        =\lim_{n\to\infty}\integral[a][b][f_n(x)][x].
    \]
\end{snippettheorem}

\begin{snippetproof}{uniform-convergence-integration-theorem-proof}{uniform-convergence-integration-theorem}{Integration of a uniformly convergent sequence}
    Write \(\lVert g\rVert_\infty=\sup_{x\in[a,b]}|g(x)|\). Fix
    \(\varepsilon>0\). Uniform convergence gives an index \(n\) such that
    \[
        \lVert f-f_n\rVert_\infty<\frac{\varepsilon}{4(b-a)}.
    \]
    Since \(f_n\) is Riemann integrable, there is a partition \(P\) of
    \([a,b]\) satisfying
    \[
        U(f_n,P)-L(f_n,P)<\frac{\varepsilon}{2},
    \]
    where \(U\) and \(L\) are the upper and lower Darboux sums. On every
    subinterval \(I\) of \(P\),
    \[
        \sup_I f-\inf_I f
        \leq \sup_I f_n-\inf_I f_n+2\lVert f-f_n\rVert_\infty.
    \]
    Multiplying by the subinterval lengths and summing gives
    \begin{align*}
        U(f,P)-L(f,P)
        &\leq U(f_n,P)-L(f_n,P)
        +2(b-a)\lVert f-f_n\rVert_\infty\\
        &<\varepsilon.
    \end{align*}
    The Darboux criterion therefore shows that \(f\) is Riemann integrable.

    Moreover,
    \begin{align*}
        \left|
            \integral[a][b][f_n(x)][x]-\integral[a][b][f(x)][x]
        \right|
        &\leq\integral[a][b][|f_n(x)-f(x)|][x]\\
        &\leq(b-a)\lVert f_n-f\rVert_\infty\longrightarrow0,
    \end{align*}
    which proves the asserted interchange.
\end{snippetproof}

\section{Unbounded domains}

\begin{plainhtml}
    Uniform convergence alone does not extend the theorem to improper integrals on unbounded domains.
\end{plainhtml}

\begin{snippetexample}{uniform-limit-improper-integrability-counterexample}{A uniform limit of improperly integrable functions need not be integrable}
    On \(S=[1,+\infty)\), define
    \[
        f_n(x)=\frac1x\chi_{[1,n]}(x)
        =\begin{cases}
            \frac1x & x\in[1,n],\\
            0 & x>n.
        \end{cases}
    \]
    For every fixed \(n\), the improper integral exists and equals
    \[
        \integral[1][\infty][f_n(x)][x]
        =\integral[1][n][\frac1x][x]=\log n.
    \]
    The sequence converges uniformly to \(f(x)=1/x\), since
    \[
        \sup_{x\in[1,+\infty)}|f_n(x)-f(x)|
        =\sup_{x>n}\frac1x=\frac1n\longrightarrow0.
    \]
    Nevertheless, \(1/x\) is not improperly integrable on
    \([1,+\infty)\).
\end{snippetexample}

\begin{snippetexample}{uniform-convergence-improper-integral-interchange-counterexample}{Uniform convergence does not justify interchanging improper integrals}
    On \(S=\realnumbers\), define
    \[
        f_n(x)=\frac1n\chi_{[n,2n]}(x)
        =\begin{cases}
            \frac1n & x\in[n,2n],\\
            0 & x\notin[n,2n].
        \end{cases}
    \]
    Every \(f_n\) is improperly integrable and
    \[
        \integral[-\infty][+\infty][f_n(x)][x]
        =\integral[n][2n][\frac1n][x]=1.
    \]
    On the other hand,
    \[
        \sup_{x\in\realnumbers}|f_n(x)|=\frac1n\longrightarrow0,
    \]
    so \(f_n\to0\) uniformly on \(\realnumbers\). Therefore
    \[
        \integral[-\infty][+\infty][0][x]
        =0\neq1
        =\lim_{n\to\infty}\integral[-\infty][+\infty][f_n(x)][x].
    \]
\end{snippetexample}

\end{document}
